\paragraph{IV 20.14: Projections coniques (cote 25, sites 84--91).}
This section supplies the systematic account of conical projection already used at the end of section 8 and in 10.4; its opening note proposes moving the account earlier, perhaps into the auxiliary Grassmannian discussion (site 84). From $C=P(E/G)$ of codimension $m+1$, Grothendieck recalls the Chapter II construction $P-C\to P(G)$ and interprets $P(G)$ as the closed family $Q(C)$ of linear spaces containing $C$. For a morphism $f:X\to P$, he then gives the restricted projection and its extension through the incidence scheme, introducing the notation $\widetilde X(C)$ and $\widetilde F(C)$ and distinguishing the latter $F$ from the quotient used at the start (sites 85--86). The generic-centre discussion recalls 5.3 as extended in section 12 and uses 10.2 for the blow-up interpretation. Proposition 14.1, substantially revised at site 87, concerns fiber dimensions; its proof at site 88 returns to section 2 as revised in section 10. Corollaries 14.2 and 14.3 treat finiteness, dominance and the Cohen--Macaulay and flatness conclusions. The added line before 14.3 refers to the remark on manuscript page 21 before 5.6 and writes $\varphi^{(m)}_\xi$, with the small subscript on the following $F$ still uncertain. Site 89 returns to the differential analysis of section 12, preserving the author's qualification ``ou montrerait, si nous ne la faisons pas'', and invokes the section 9 replacement by a Veronese multiple. Proposition 14.4 concerns a finite subscheme $T$: the revised item a) adds the immersion hypothesis and replaces ``un monomorphisme'' by ``radiciel i.e. géométriquement injectif''; item b) receives a further handwritten inverse-image clause. Proposition 14.5 and its added conclusion lead to Corollary 14.6 at site 90, where the original birational-isomorphism wording is replaced by an existential formulation involving a ``plongement birationnel''. The following discussion records limits to the proposed immersion and nonramification conclusions, explicitly mentions Chapter VIII, and explains the varying ramification locus. Site 91 closes with three base-change diagrams and the universal family, returning to 7.1 and proposing $Q$ as the name of a morphism of functors; ``Hilb'' is canceled in its target. The revisions should be read alongside their scans: at site 86 the generic centre is literally called a linear subvariety of $X_{k(\eta)}$; in the amended 14.1(b) at site 87 the deletion crosses the printed ``si'' together with the old properness clause; and the dimension condition at site 89 is a typed interlinear insertion whose connector places it after the parenthetical qualification. These page readings and the remaining shorthand doubts are retained without supplying a revised mathematical statement.
